% =============================================================================
% ---> Геометрия страницы
% =============================================================================

\usepackage[
  papersize={170mm,240mm},
  inner=18mm,
  outer=20mm,
  bindingoffset=10mm,
  top=25mm,
  bottom=28mm,
  headheight=0pt,
  headsep=0pt,
  footskip=13mm,
  heightrounded,
]{geometry}

\raggedbottom


% =============================================================================
% ---> Шрифты
% =============================================================================

\usepackage{fontspec}

\defaultfontfeatures{Ligatures=TeX}

\setmainfont{STIX Two Text}[
  BoldFont       = STIX Two Text Bold,
  ItalicFont     = STIX Two Text Italic,
  BoldItalicFont = STIX Two Text Bold Italic,
  Script         = Default,
  Language       = Default,
]

\setsansfont{Libertinus Sans}[
  BoldFont       = Libertinus Sans Bold,
  ItalicFont     = Libertinus Sans Italic,
  BoldItalicFont = Libertinus Sans Italic, % Bold Italic не существует!
  Script         = Default,
  Language       = Default,
]

\directlua{
  luaotfload.add_fallback("monofallback", {
    "DejaVu Sans Mono:mode=node;",
    % при желании можно добавить ещё:
    % "STIX Two Math:mode=node;",
  })
}

\setmonofont{CMU Typewriter Text}[
  BoldFont       = CMU Typewriter Text Bold,
  ItalicFont     = CMU Typewriter Text Italic,
  BoldItalicFont = CMU Typewriter Text BoldItalic,
  Script         = Default,
  Language       = Default,
  RawFeature     = {fallback=monofallback},
]

\DeclareFontSeriesDefault[sf]{bf}{b}


% =============================================================================
% ---> Математика
% =============================================================================

\usepackage{amsmath}
\usepackage{amsthm}
\usepackage{mathtools}
\usepackage{unicode-math}

\setmathfont{STIX Two Math}


% =============================================================================
% ---> Цветные блоки и рамки
% =============================================================================

\usepackage[most]{tcolorbox}


% =============================================================================
% ---> Язык
% =============================================================================

\usepackage{polyglossia}

\setdefaultlanguage{russian}
\setotherlanguage{english}


% =============================================================================
% ---> Графика и рисунки
% =============================================================================

\usepackage{graphicx}
\usepackage{caption}
\usepackage{subcaption}

\captionsetup{
  font=small,
  labelfont=bf,
  justification=raggedright,
  singlelinecheck=false,
}


% =============================================================================
% ---> Форматирование текста
% =============================================================================

\usepackage{indentfirst}
\usepackage{enumitem}
\usepackage{lua-ul}
\usepackage{etoolbox}
\usepackage{fontawesome}
\usepackage{xspace}
\usepackage{longtable}
\usepackage{array}
\usepackage{lastpage}


% =============================================================================
% ---> TikZ
% =============================================================================

\usepackage{tikz}

\usetikzlibrary{
  calc,
  decorations.pathmorphing,
  fadings,
}


% =============================================================================
% ---> Код на Python (minted)
% =============================================================================

\usepackage{minted}

\setminted{
  fontsize=\small,
  breaklines=true,
  breakanywhere=false,
  autogobble=true,
  tabsize=4,
  frame=none,
  numbersep=8pt,
  xleftmargin=0pt,
  xrightmargin=0pt,
}
\setminted[python]{linenos=false}


% =============================================================================
% ---> Псевдокод алгоритмов (algpseudocode)
% =============================================================================

\usepackage{algpseudocode}

% Русификация ключевых слов
\algrenewcommand\algorithmicif{\textbf{если}}
\algrenewcommand\algorithmicthen{\textbf{то}}
\algrenewcommand\algorithmicelse{\textbf{иначе}}
\algrenewcommand\algorithmicend{\textbf{конец}}
\algrenewcommand\algorithmicfor{\textbf{для}}
\algrenewcommand\algorithmicforall{\textbf{для всех}}
\algrenewcommand\algorithmicdo{\textbf{выполнять}}
\algrenewcommand\algorithmicwhile{\textbf{пока}}
\algrenewcommand\algorithmicrepeat{\textbf{повторять}}
\algrenewcommand\algorithmicuntil{\textbf{пока не}}
\algrenewcommand\algorithmicreturn{\textbf{вернуть}}
\algrenewcommand\algorithmicrequire{\textbf{Вход:}}
\algrenewcommand\algorithmicensure{\textbf{Выход:}}


% =============================================================================
% ---> Списки
% =============================================================================

\setlist[itemize,1]{
  label=\textcolor{listL1}{\textbullet},
  leftmargin=\parindent,
  itemsep=0.0em,
  parsep=0pt,
  topsep=0.5em,
  partopsep=0.5em,
}
\setlist[itemize,2]{
  label=\textcolor{listL2}{\textbullet},
  leftmargin=*,
  itemsep=0.25em,
  parsep=0pt,
  topsep=0.5em,
}
\setlist[itemize,3]{
  label=\textcolor{listL3}{\textbullet},
  leftmargin=*,
  itemsep=0.2em,
  parsep=0pt,
  topsep=0.4em,
}

\setlist[enumerate,1]{
  label=\textcolor{listL1}{\arabic*.},
  leftmargin=\parindent,
  itemsep=0pt,
  parsep=0pt,
  topsep=0.5em,
  partopsep=0.5em,
}
\setlist[enumerate,2]{
  label=\textcolor{listL2}{\arabic{enumi}.\arabic*.},
  leftmargin=*,
  itemsep=0.35em,
  parsep=0pt,
  topsep=0.5em,
}
\setlist[enumerate,3]{
  label=\textcolor{listL3}{\arabic{enumi}.\arabic{enumii}.\arabic*.},
  leftmargin=*,
  itemsep=0.35em,
  parsep=0pt,
}


% =============================================================================
% ---> Колонтитулы
% =============================================================================

\usepackage{fancyhdr}

\newcommand{\mypagenum}{%
  \textcolor{xgray}{\textbf{\textemdash}} \hspace{0.5em}%
  \textbf{\thepage}%
  \hspace{0.5em} \textcolor{xgray}{\textbf{\textemdash}}%
}

\pagestyle{fancy}
\fancyhf{}
\fancyfoot[C]{\mypagenum}
\renewcommand{\headrulewidth}{0pt}
\renewcommand{\footrulewidth}{0pt}

\fancypagestyle{plain}{
  \fancyhf{}
  \fancyfoot[C]{\mypagenum}
  \renewcommand{\headrulewidth}{0pt}
  \renewcommand{\footrulewidth}{0pt}
}


% =============================================================================
% ---> Заголовки глав, секций и подсекций
% =============================================================================

\usepackage{titlesec}

% =============================================================================
% style/base_chapter.tex
% =============================================================================


% =============================================================================
% Двузначный номер главы в бейдже; в приложениях оставляем A, B, ...
% =============================================================================

\makeatletter
\@ifundefined{ifchapterbadgeappendix}{\newif\ifchapterbadgeappendix}{}
\makeatother

\pretocmd{\appendix}{\chapterbadgeappendixtrue}{}{}

\providecommand{\chapterBadgeNumber}{}
\renewcommand{\chapterBadgeNumber}{%
  \ifchapterbadgeappendix%
    \thechapter%
  \else%
    \ifnum\value{chapter}<10 0\fi\arabic{chapter}%
  \fi%
}


% =============================================================================
% Цвета заголовков глав
% =============================================================================

% Нумерованный шестиугольный бейдж.
\colorlet{chapterBadgeTop}{xmatrix!55!white}
\colorlet{chapterBadgeBottom}{xmatrix!85!black}
\colorlet{chapterBadgeFrame}{xmatrix!45!black}
\colorlet{chapterBadgeText}{xmatrix!3!white}

% Одна горизонтальная линия под нумерованной главой.
\colorlet{chapterRuleColor}{xmatrix!60!black}

% Толстая вертикальная метка для ненумерованной главы.
\colorlet{chapterNumberlessLine}{xmatrix!60!black}


% =============================================================================
% Fading для чистого верхнего блика
% =============================================================================

\tikzfading[
  name=chapterGlossFade,
  top color=transparent!0,
  bottom color=transparent!100,
]


% =============================================================================
% Геометрия и выравнивание декоративных элементов заголовка
% =============================================================================

\newlength{\chapterBadgeBaselineShift}
\setlength{\chapterBadgeBaselineShift}{6.3pt}

\newlength{\chapterNumberlessBaselineShift}
\setlength{\chapterNumberlessBaselineShift}{6.3pt}

% Контур шестиугольного бейджа.
\newcommand{\chapterBadgeHexPath}[1]{%
  ([xshift=2.8mm]#1.north west) --
  ([xshift=-2.8mm]#1.north east) --
  ([yshift=-0.25mm]#1.east) --
  ([xshift=-2.8mm]#1.south east) --
  ([xshift=2.8mm]#1.south west) --
  ([yshift=0.25mm]#1.west) --
  cycle%
}


% =============================================================================
% Бейдж нумерованной главы
% =============================================================================

\providecommand{\chapterMatrixBadgeE}[1]{}
\renewcommand{\chapterMatrixBadgeE}[1]{%
  \tikz[baseline={([yshift=-\chapterBadgeBaselineShift]B.center)}]{%
    % Невидимый узел задаёт размеры бейджа.
    \node[
      minimum width=18.4mm,
      minimum height=12.0mm,
      inner xsep=5.8pt,
      inner ysep=1.9pt,
      outer sep=0pt,
      font=\sffamily\bfseries\fontsize{20}{22}\selectfont
    ] (B) {\phantom{\strut #1}};

    % Внешняя мягкая тень.
    \begin{scope}[xshift=0.50mm, yshift=-0.60mm]
      \path[
        fill=black,
        opacity=0.18,
        rounded corners=1.8pt,
      ]
        \chapterBadgeHexPath{B};
    \end{scope}

    % Основная заливка: чистый вертикальный градиент без среднего цвета.
    \shade[
      top color=chapterBadgeTop,
      bottom color=chapterBadgeBottom,
      draw=chapterBadgeFrame,
      line width=0.80pt,
      rounded corners=1.8pt,
      line join=round,
    ]
      \chapterBadgeHexPath{B};

    % Верхний глянцевый блик: плавно исчезает к середине бейджа.
    \begin{scope}
      \clip[rounded corners=1.8pt]
        \chapterBadgeHexPath{B};
      \path[
        fill=white,
        path fading=chapterGlossFade,
        opacity=0.55,
      ]
        ([yshift=-0.25mm]B.north west) rectangle
        ([yshift=-0.25mm]B.east);
    \end{scope}

    % Тонкий внутренний световой кант вверху.
    \begin{scope}
      \clip[rounded corners=1.8pt]
        \chapterBadgeHexPath{B};
      \draw[
        white,
        opacity=0.55,
        line width=0.50pt,
        line cap=round,
      ]
        ([xshift=3.0mm, yshift=-0.42mm]B.north west) --
        ([xshift=-3.0mm, yshift=-0.42mm]B.north east);
    \end{scope}

    % Тонкий внутренний теневой кант внизу.
    \begin{scope}
      \clip[rounded corners=1.8pt]
        \chapterBadgeHexPath{B};
      \draw[
        black,
        opacity=0.22,
        line width=0.45pt,
        line cap=round,
      ]
        ([xshift=3.0mm, yshift=0.45mm]B.south west) --
        ([xshift=-3.0mm, yshift=0.45mm]B.south east);
    \end{scope}

    % Наружная обводка поверх всех эффектов.
    \path[
      draw=chapterBadgeFrame,
      line width=0.80pt,
      rounded corners=1.8pt,
      line join=round,
    ]
      \chapterBadgeHexPath{B};

    % Номер главы.
    \node[
      minimum width=18.4mm,
      minimum height=12.0mm,
      inner xsep=5.8pt,
      inner ysep=1.9pt,
      outer sep=0pt,
      text=chapterBadgeText,
      font=\sffamily\bfseries\fontsize{20}{22}\selectfont,
    ] at (B.center) {\strut #1};
  }%
}


% =============================================================================
% Метка ненумерованной главы: одна толстая вертикальная линия
% =============================================================================

\providecommand{\chapterNumberlessMark}{}
\renewcommand{\chapterNumberlessMark}{%
  \tikz[baseline={([yshift=-\chapterNumberlessBaselineShift]C.center)}]{%
    \node[
      minimum width=5.0mm,
      minimum height=12.0mm,
      inner sep=0pt,
      outer sep=0pt
    ] (C) {};

    \draw[
      chapterNumberlessLine,
      line width=2.6pt,
      line cap=round,
    ]
      ([xshift=2.2mm, yshift=-0.3mm]C.north west) --
      ([xshift=2.2mm, yshift=0.3mm]C.south west);
  }%
}


% =============================================================================
% Нумерованные главы: бейдж + одна толстая горизонтальная линия
% =============================================================================

\titleformat{\chapter}[hang]
  {%
    \normalfont\sffamily\bfseries\raggedright\color{xmatrixDark}%
    \hyphenpenalty=10000\exhyphenpenalty=10000%
  }
  {\chapterMatrixBadgeE{\chapterBadgeNumber}}
  {13pt}
  {\fontsize{24}{28}\selectfont}
  [\vspace{7.5pt}{\color{chapterRuleColor}\titlerule[1.90pt]}]


% =============================================================================
% Ненумерованные главы: только вертикальная метка, без горизонтальной линии
% =============================================================================

\titleformat{name=\chapter,numberless}[hang]
  {%
    \normalfont\sffamily\bfseries\raggedright\color{xmatrixDark}%
    \hyphenpenalty=10000\exhyphenpenalty=10000%
  }
  {\chapterNumberlessMark}
  {13pt}
  {\fontsize{24}{28}\selectfont}


% =============================================================================
% Отступы заголовка главы
% =============================================================================

\titlespacing*{\chapter}{0pt}{10pt plus 2pt minus 2pt}{20pt}    % Стиль оформления chapter
% =============================================================================
% style/base_section.tex
% =============================================================================


% =============================================================================
% Цвета заголовков разделов
% =============================================================================

%\colorlet{sectionAccent}{xindigo!78!xblue}
\colorlet{sectionAccent}{xindigo!55!black}
\colorlet{sectionAccentDark}{sectionAccent!55!black}

% Цвета номерного бейджа.
\colorlet{sectionBadgeTop}{sectionAccent!72!black}
\colorlet{sectionBadgeBottom}{sectionAccent!28!black}
\colorlet{sectionBadgeFrame}{sectionAccent!38!black}
\colorlet{sectionBadgeText}{sectionAccent!6!white}

% Цвет заголовка и двух нижних линий.
\colorlet{sectionTitleColor}{sectionAccentDark}
\colorlet{sectionRuleA}{sectionAccent!65!white}
\colorlet{sectionRuleB}{sectionAccent!45!white}

% Две вертикальные линии для ненумерованного section.
\colorlet{sectionNumberlessLineA}{sectionAccent!65!white}
\colorlet{sectionNumberlessLineB}{sectionAccent!75!black}


% =============================================================================
% Fading для чистого верхнего блика
% =============================================================================

\tikzfading[
  name=sectionGlossFade,
  top color=transparent!0,
  bottom color=transparent!100,
]


% =============================================================================
% Геометрия
% =============================================================================

\newlength{\sectionBadgeBaselineShift}
\setlength{\sectionBadgeBaselineShift}{5.2pt}

\newlength{\sectionNumberlessBaselineShift}
\setlength{\sectionNumberlessBaselineShift}{5.2pt}


% =============================================================================
% Номерной глянцевый бейдж для section
% =============================================================================

\providecommand{\sectionGlossyBadge}[1]{}
\renewcommand{\sectionGlossyBadge}[1]{%
  \tikz[baseline={([yshift=-\sectionBadgeBaselineShift]S.center)}]{%
    % Невидимый узел задаёт размеры бейджа.
    \node[
      minimum width=13.6mm,
      minimum height=7.8mm,
      inner xsep=4.4pt,
      inner ysep=1.3pt,
      outer sep=0pt,
      font=\sffamily\bfseries\fontsize{10.8}{12}\selectfont
    ] (S) {\phantom{\strut #1}};

    % Внешняя тень.
    \begin{scope}[xshift=0.38mm, yshift=-0.46mm]
      \path[
        fill=black,
        opacity=0.16,
        rounded corners=2.0pt,
      ]
        (S.south west) rectangle (S.north east);
    \end{scope}

    % Основная заливка: чистый вертикальный градиент.
    \shade[
      top color=sectionBadgeTop,
      bottom color=sectionBadgeBottom,
      draw=sectionBadgeFrame,
      line width=0.58pt,
      rounded corners=2.0pt,
    ]
      (S.south west) rectangle (S.north east);

    % Верхний глянцевый блик: плавно исчезает к середине бейджа.
    \begin{scope}
      \clip[rounded corners=2.0pt]
        (S.south west) rectangle (S.north east);
      \path[
        fill=white,
        path fading=sectionGlossFade,
        opacity=0.52,
      ]
        ([yshift=-0.18mm]S.north west) rectangle
        ([yshift=-0.18mm]S.east);
    \end{scope}

    % Верхний внутренний световой кант.
    \begin{scope}
      \clip[rounded corners=2.0pt]
        (S.south west) rectangle (S.north east);
      \draw[
        white,
        opacity=0.55,
        line width=0.40pt,
        line cap=round,
      ]
        ([xshift=1.1mm, yshift=-0.38mm]S.north west) --
        ([xshift=-1.1mm, yshift=-0.38mm]S.north east);
    \end{scope}

    % Нижний внутренний теневой кант.
    \begin{scope}
      \clip[rounded corners=2.0pt]
        (S.south west) rectangle (S.north east);
      \draw[
        black,
        opacity=0.20,
        line width=0.38pt,
        line cap=round,
      ]
        ([xshift=1.1mm, yshift=0.40mm]S.south west) --
        ([xshift=-1.1mm, yshift=0.40mm]S.south east);
    \end{scope}

    % Обводка поверх бликов.
    \path[
      draw=sectionBadgeFrame,
      line width=0.58pt,
      rounded corners=2.0pt,
    ]
      (S.south west) rectangle (S.north east);

    % Номер section.
    \node[
      minimum width=13.6mm,
      minimum height=7.8mm,
      inner xsep=4.4pt,
      inner ysep=1.3pt,
      outer sep=0pt,
      text=sectionBadgeText,
      font=\sffamily\bfseries\fontsize{10.8}{12}\selectfont,
    ] at (S.center) {\strut #1};
  }%
}


% =============================================================================
% Метка ненумерованного section: две вертикальные линии
% =============================================================================

\providecommand{\sectionNumberlessMark}{}
\renewcommand{\sectionNumberlessMark}{%
  \tikz[baseline={([yshift=-\sectionNumberlessBaselineShift]S.center)}]{%
    \node[
      minimum width=5.0mm,
      minimum height=7.8mm,
      inner sep=0pt,
      outer sep=0pt
    ] (S) {};

    \draw[
      sectionNumberlessLineA,
      line width=1.30pt,
      line cap=round,
    ]
      ([xshift=1.4mm, yshift=-0.2mm]S.north west) --
      ([xshift=1.4mm, yshift=0.2mm]S.south west);

    \draw[
      sectionNumberlessLineB,
      line width=1.50pt,
      line cap=round,
    ]
      ([xshift=3.1mm, yshift=-0.2mm]S.north west) --
      ([xshift=3.1mm, yshift=0.2mm]S.south west);
  }%
}


% =============================================================================
% Две линии под нумерованным section
% =============================================================================

\providecommand{\sectionDoubleRule}{}
\renewcommand{\sectionDoubleRule}{%
  \vspace{3.0pt}%
  {\color{sectionRuleA}\titlerule[1.00pt]}%
  \vspace{1.25pt}%
  {\color{sectionRuleB}\titlerule[0.65pt]}%
}


% =============================================================================
% Нумерованные section: бейдж + две горизонтальные линии
% =============================================================================

\titleformat{\section}[hang]
  {%
    \normalfont\sffamily\bfseries\raggedright\color{sectionTitleColor}%
    \hyphenpenalty=10000\exhyphenpenalty=10000%
  }
  {\sectionGlossyBadge{\thesection}}
  {10pt}
  {\fontsize{18}{22}\selectfont}
  [\sectionDoubleRule]


% =============================================================================
% Ненумерованные section: только вертикальная метка, без горизонтальных линий
% =============================================================================

\titleformat{name=\section,numberless}[hang]
  {%
    \normalfont\sffamily\bfseries\raggedright\color{sectionTitleColor}%
    \hyphenpenalty=10000\exhyphenpenalty=10000%
  }
  {\sectionNumberlessMark}
  {10pt}
  {\fontsize{18}{22}\selectfont}


% =============================================================================
% Отступы section
% =============================================================================

\titlespacing*{\section}
  {0pt}
  {5.7ex plus 1.8ex minus 1.4ex}
  {2.6ex plus 0.6ex minus 0.5ex}    % Стиль оформления section
% =============================================================================
% style/base_subsection.tex
% =============================================================================


% =============================================================================
% Цвета заголовков подразделов
% =============================================================================

\colorlet{subsectionAccent}{xred!72!black}
\colorlet{subsectionTitleColor}{subsectionAccent}

% Три тонкие вертикальные линии (subsection всегда без номера).
\colorlet{subsectionMarkA}{xred!60!white}
\colorlet{subsectionMarkB}{xred!75!white}
\colorlet{subsectionMarkC}{xred!90!black}


% =============================================================================
% Геометрия
% =============================================================================

\newlength{\subsectionBaselineShift}
\setlength{\subsectionBaselineShift}{4.6pt}


% =============================================================================
% Метка subsection: три тонкие вертикальные линии
% =============================================================================

\providecommand{\subsectionMark}{}
\renewcommand{\subsectionMark}{%
  \tikz[baseline={([yshift=-\subsectionBaselineShift]S.center)}]{%
    \node[
      minimum width=5.2mm,
      minimum height=5.4mm,
      inner sep=0pt,
      outer sep=0pt
    ] (S) {};

    \draw[
      subsectionMarkA,
      line width=0.95pt,
      line cap=round,
    ]
      ([xshift=0.8mm, yshift=-0.15mm]S.north west) --
      ([xshift=0.8mm, yshift=0.15mm]S.south west);

    \draw[
      subsectionMarkB,
      line width=1.10pt,
      line cap=round,
    ]
      ([xshift=2.2mm, yshift=-0.15mm]S.north west) --
      ([xshift=2.2mm, yshift=0.15mm]S.south west);

    \draw[
      subsectionMarkC,
      line width=1.25pt,
      line cap=round,
    ]
      ([xshift=3.6mm, yshift=-0.15mm]S.north west) --
      ([xshift=3.6mm, yshift=0.15mm]S.south west);
  }%
}


% =============================================================================
% Subsection: метка слева, без горизонтальных линий
% =============================================================================

\titleformat{\subsection}[hang]
  {%
    \normalfont\sffamily\bfseries\fontsize{15.5}{19}\selectfont%
    \raggedright\color{subsectionTitleColor}%
    \hyphenpenalty=10000\exhyphenpenalty=10000%
  }
  {\subsectionMark}
  {0.55em}
  {}


% =============================================================================
% Звёздная форма subsection
% =============================================================================

\titleformat{name=\subsection,numberless}[hang]
  {%
    \normalfont\sffamily\bfseries\fontsize{15.5}{19}\selectfont%
    \raggedright\color{subsectionTitleColor}%
    \hyphenpenalty=10000\exhyphenpenalty=10000%
  }
  {\subsectionMark}
  {0.55em}
  {}


% =============================================================================
% Отступы subsection
% =============================================================================

\titlespacing*{\subsection}
  {0pt}
  {4.5ex plus 1.5ex minus 1.3ex}
  {2.0ex plus 0.5ex minus 0.4ex} % Стиль оформления subsection


% =============================================================================
% ---> Микротипографика
% =============================================================================

\usepackage{microtype}

\emergencystretch=3em

\clubpenalty=10000
\widowpenalty=10000
\displaywidowpenalty=10000
\brokenpenalty=10000
\binoppenalty=4000
\relpenalty=4000

\AtBeginDocument{%
  \setlength{\abovedisplayskip}{8pt plus 3pt minus 3pt}%
  \setlength{\belowdisplayskip}{8pt plus 3pt minus 3pt}%
  \setlength{\abovedisplayshortskip}{2pt plus 3pt minus 2pt}%
  \setlength{\belowdisplayshortskip}{5pt plus 3pt minus 2pt}%
}
\allowdisplaybreaks[1]

\setlength{\skip\footins}{18pt plus 4pt minus 2pt}

\renewcommand{\footnoterule}{%
  \kern-3pt
  \hrule width 0.35\textwidth height 0.4pt
  \kern 2.6pt
}


% =============================================================================
% ---> Гиперссылки
% =============================================================================

\usepackage[bookmarks, hidelinks]{hyperref}
\usepackage{bookmark}
\hypersetup{
  pdftitle          = {\PDFBookTitle},
  pdfauthor         = {\PDFBookAuthor},
  pdfsubject        = {\PDFBookSubject},
  pdfkeywords       = {\PDFBookKeywords},
  pdflang           = {ru-RU},
  pdfencoding       = unicode,
  bookmarksnumbered = true,
  bookmarksopen     = true,
}
\pdfstringdefDisableCommands{%
  \def\foc#1{#1}%
  \def\emp#1{#1}%
  \def\underLine#1{#1}%
  \def\textcolor#1#2{#2}%
}
\urlstyle{same}


% =============================================================================
% ---> Счётчики и нумерация
% =============================================================================

\setcounter{tocdepth}{1}
\setcounter{secnumdepth}{1}


% =============================================================================
% ---> Локализация
% =============================================================================

\addto\captionsrussian{
  %\renewcommand{\chaptername}{Лекция}
  \renewcommand{\contentsname}{Содержание}
}


% =============================================================================
% ---> Заметки и задачи (todonotes)
% =============================================================================
% Для отключения добавить disable: \usepackage[disable, ...]{todonotes}

\usepackage[textsize=scriptsize, colorinlistoftodos, shadow, loadshadowlibrary]{todonotes}
\makeatletter
\providecommand*{\toclevel@todo}{0}
\makeatother